\begin{definition}
An \emph{event structure} is a tuple $(E, {\leq}, \#)$ where $\leq$ is a partial order on $E$, and $\#$  is an irreflexive and symmetric binary relation on $E$, satisfying the axioms below:
\begin{itemize}
    \item For every $e \in E$  the set $\{e' \in E \mid e' \leq e \}$ is finite.
    \item For $e, e', e'' \in E$, if $e \leq_{E} e'$ and $e \mathrel{\#} e''$ then $e' \mathrel{\#} e''$.
\end{itemize}
\end{definition}
The execution paths of an event structure are called \emph{configurations}:
\begin{definition}
    A finite subset $X \subseteq E$ is called \emph{consistent} if no two events in $X$ are in conflict. It is called a \emph{configuration} if it is additionally down-closed. The set of consistent subsets of $E$ is denoted $\Con(E)$, and the set of configurations is called $\conf(E)$.
\end{definition}
Maps of event structures must faithfully preserve configurations:
\begin{definition}
    If $(E, \leq_E, \#_E)$ and $(E', \leq_{E'}, \#_{E'})$ are event structures, a function $f: E \rightarrow E'$ is a \emph{map of event structures} if for every $x \in \conf(E)$, the direct image $fx := \{f(e) \mid e \in x\}$ is in $\conf(E')$, and moreover the restriction of $f$ to $x$ is injective.
\end{definition}
Event structures and maps of event structures form a category $\ES$. Frequently, we will regard $\ES$ as a locally discrete $2$-category, just like we did for $\Petri$ and $\Occ$.

\subsection{The unfolding of a Petri net as an event structure with symmetry}

We first recall event structures with symmetry and the $2$-category $\SES$, originally introduced in [32].

\subsubsection{Event structures with symmetry}

\label{subsec:sem-def}
 Informally, symmetry on an event structure indicates which pairs of executions can be considered equivalent, via a bisimulation relation. 
\begin{definition}
    For an event structure $E$, an isomorphism family  is a set $\mathbb{S}$ of bijections $\theta:x \cong y$ between finite configurations of $x, y \in \conf(E)$, such that:
    \begin{itemize}
        \item $\mathbb{S}$ contains all identity bijections, and is stable under composition and inverse.  
        \item For $\theta : x \cong y \in \mathbb{S}$, if $x' \subseteq x$, then the restriction of $\theta$ to $x'$ is in $\mathbb{S}$.
        \item For $\theta : x \cong y \in \mathbb{S}$, if $x \subseteq x'$, then there exists an extension $y \subseteq y'$ and a bijection $\theta' : x' \cong y' \in \mathbb{S}$ such that $\theta'$ restricts to $\theta$.
    \end{itemize}
    The pair $(E, \mathbb{S})$ is called an \emph{event structure with symmetry}.
\end{definition}

We then consider symmetry-preserving maps:
\begin{definition}
    For event structures with symmetry $(E,\mathbb{S})$ and $(E', \mathbb{S}')$, a map $f : E \rightarrow E'$ \emph{preserves symmetry} if for all $\theta : x \cong y \in \mathbb{S}$, $f \theta : fx \cong fy \in \mathbb{S}'$, where $f\theta : f(e) \mapsto f(\theta(e))$.
\end{definition}
A novelty in the presence of symmetry is the following equivalence relation on maps:
\begin{definition}
    For event structures with symmetry $(E,\mathbb{S})$ and $(E', \mathbb{S}')$ and symmetry-preserving maps $f, g : E \rightarrow E'$, say that $f$ and $g$ are \emph{symmetric}, denoted $f \sim g$, if for every $x \in \conf(E)$ the bijection $fx \cong gx$ defined by $ f(e) \mapsto g(e)$ is in $\mathbb{S}'$.
\end{definition}
Together, event structures with symmetry, symmetry-preserving maps and symmetries of maps form a $2$-category $\SES$ [32]. This is a degenerate 2-category in the sense that hom-categories are all \emph{setoids} (sets with an equivalence relation). There is an embedding $J : \ES \to \SES$. As we will see, the 2-categorical structure makes it possible to consider universal properties ``up to symmetry''.

\subsection{A 2-density result for event structures}
\label{subsec:2density}

\newcommand{\bE}{\mathbf{E}}
\newcommand{\bC}{\mathbf{C}}
\newcommand{\bD}{\mathbf{D}}

We prove a density result for the embedding of event structures into event structures with symmetry. This is a 2-categorical (i.e. $\Cat$-enriched) form of density: each event structure with symmetry is a canonical 2-colimit of  ordinary event structures. Another way of stating this is as follows:

\newcommand{\op}{\mathrm{op}}

\begin{definition}[e.g. [17]]
For locally small 2-categories $\bD$ and $\bE$, a $2$-functor $J : \bD \rightarrow \bE$ is \emph{2-dense} if the 2-functor $\tilde{J} : \bE \to [\bD^\op, \Cat] : e \mapsto \bE(J-, e)$ is locally an isomorphism of categories, where $[\bD^\op, \Cat]$ is the 2-category of 2-functors, strict natural transformations, and modifications. 
\end{definition}

\begin{theorem}
The embedding 2-functor $J : \ES \to \SES$ is 2-dense.
\end{theorem}

(We emphasize that 2-density is strictly weaker than density in the 1-categorical sense. In particular the underlying 1-functor $J : \ES \to \SES$ is not dense.)
\begin{proof}
For event structures with symmetry $A$ and $B$, we show that the functor $\tilde{J}_{A ,B} : \SES(A, B) \to \Cat^{\ES^\op}(\SES(J-, A), \SES(J-, B))$ is an isomorphism of categories by constructing an inverse. Write $A = (|A|, \mathbb{S}_A)$ and $B = (|B|, \mathbb{S}_B)$. 

Observe that, for a natural transformation $\alpha : \SES(J-, A) \to \SES(J-, B)$, applying component $\alpha_{|A|}$ to the `identity' function $\id_{|A|} : J|A| \to A$ gives a map of event structures with symmetry $\alpha_{|A|}(\id_{|A|}) : J|A| \to B$. We claim that $\alpha_{|A|}(\id_{|A|})$ preserves the symmetry in $A$. For any bijection $\theta: x \cong y$ in $\mathbb{S}_A$, we have $(\id_{|A|} \circ J\iota_y \circ J\theta) \sim (\id_{|A|} \circ J\iota_x)$, writing $\iota_x : x \to |A|$ and $\iota_y : y \to |A|$ for the inclusion maps (and recalling that $\theta \in \ES(x, y)$). Since $\alpha$ is a natural transformation and each component preserves the symmetry relation on maps,  $\alpha_{|A|}(\id_{|A|}) \circ J{\iota_x} = \alpha_x (\id_{|A|} \circ J{\iota_x}) \sim \alpha_x(\id_{|A|} \circ J{\iota_y} \circ J\theta) = \alpha_{|A|}(\id_{|A|}) \circ J\iota_y \circ J\theta.$ In other words, $\alpha_{|A|}(\id_{|A|}) \theta \in \mathbb{S}_B$, and we have proven the claim. Therefore, we take $\tilde J_{A, B}^{-1}(\alpha) = \alpha_{|A|}(\id_{|A|})$. For functoriality of $J_{A, B}^{-1}$, note that any modification $\alpha \to \beta$ must be unique and gives $\alpha_{|A|}(\id_{|A|})\sim \beta_{|A|}(\id_{|A|})$. We omit the verification that $\tilde J_{A, B}^{-1}$ and $\tilde J_{A, B}$ are inverses. 
\end{proof}

